\section{Emergence of Machine Learning in Materials Science} \label{section:machine_learning}

% ---------- %

Over the past decade, machine learning (ML) has emerged as a transformative tool in materials modelling, offering a
route to accelerate predictions and expand accessible system sizes beyond those tractable by first-principles methods.

Traditional computational techniques, such as DFT and MD, have established themselves as accurate and rigorous, but they
often scale poorly with system size or configurational complexity.

This limitation becomes particularly acute in the modelling of interfaces, where subtle geometric, chemical, and
electronic variations can dramatically affect material behaviour.

ML offers a complementary paradigm, enabling the construction of surrogate models trained on high-fidelity data, which
can interpolate or even extrapolate to predict key material properties across vast configuration spaces.

% ---------- %

This section introduces the growing role of ML in materials prediction, highlighting its application in interface
modelling and the broader landscape of data-driven approaches.

Subsequent subsections will examine the types of models in use, including classical regressors and deep neural
architectures, alongside the nature of training data, typical input descriptors, and validation strategies.

Special attention will be given to machine-learned interatomic potentials (MLPs), particularly those with demonstrated
success in interfacial contexts such as MACE and ChgNet.

Finally, we reflect on the benefits and ongoing limitations of ML methods in materials science, including issues of
transferability, uncertainty quantification, and integration into automated modelling pipelines.

% ---------- %

\subsection{How has machine learning been introduced into materials prediction?}

% ---------- %

Machine learning (ML) has been introduced into materials prediction as a complementary tool to traditional
first-principles methods, offering improved scalability, predictive flexibility, and the potential for workflow
automation.

Its adoption has been driven by limitations in conventional approaches such as Density Functional Theory (DFT),
particularly in capturing the structural and chemical complexity of real-world systems such as interfaces.

% ---------- %

ML enables efficient exploration of vast configurational landscapes, where brute-force DFT sampling would be infeasible.

This is particularly important in interface modelling, where multiple stacking orders, terminations, and slip
configurations can exist.

ML-driven screening allows researchers to rapidly identify promising candidates for further evaluation.

% ---------- %

Early ML applications focused on surrogate models trained to replicate DFT-derived energies, forces, or properties.

Descriptor-based regression models, using quantities such as coordination numbers, Bader charges, or bond lengths, were
employed to predict interfacial stability or band alignments.

These approaches have evolved toward end-to-end learning frameworks using physically-informed descriptors or graph-based
representations.

% ---------- %

ML is now embedded into iterative prediction--generation--evaluation cycles.

For example, predicted interface favourability informs structure generation (e.g. stacking vectors or surface
terminations), which are then evaluated by DFT or ML surrogates.

The results feed back into model training, supporting a self-refining loop.

Tools such as ARTEMIS benefit from such integration.

% ---------- %

More recent work employs graph neural networks (GNNs), such as MACE or NequIP, which are equivariant to Euclidean
transformations and thus well-suited for atomistic systems.

These models learn directly from atomic positions and species, avoiding the need for handcrafted features while
preserving symmetry constraints essential for interface modelling.

% ---------- %

Ultimately, ML enables a shift from static modelling pipelines to more autonomous workflows.

When coupled with structure-generation protocols and physical constraints.

% ---------- %

\subsection{What types of models and data are being used?}

% ---------- %

Models can be broadly categorised into classical learning algorithms and deep learning approaches.

Classical methods include random forests, gradient boosting, support vector machines (SVMs), and Gaussian process
regression (GPR), the latter being particularly valuable for small datasets due to its uncertainty estimation
capability.

Deep learning architectures include feedforward neural networks and, more notably, graph neural networks (GNNs), which
are well-suited for materials modelling due to their ability to represent variable-sized atomic graphs with complex
bonding topologies.

% ---------- %

The input to these models typically consists of numerical representations of atomic configurations known as descriptors.

Geometric descriptors include bond lengths, coordination numbers, angular distributions, $n$-body distribution
functions, and radial symmetry functions; many of which feature prominently in Behler-Parrinello-type neural networks.

Chemical descriptors such as electronegativity differences and Bader charges capture bonding character and charge
transfer.

Electronic descriptors like partial density of states (PDOS) and local work functions are relevant for predicting band
alignment and interfacial transport properties.

Interface-specific features, such as interfacial roughness or vacuum level offsets, are particularly crucial in 2D|3D
and mixed-dimensional heterostructures.

% ---------- %

To support learning across such heterogeneous and high-dimensional features, symmetry-invariant and scale-independent
representations have become increasingly important.

These ensure that models generalise across materials with different atom counts or lattice sizes.

Dimensionality reduction methods and learned embeddings (e.g. via autoencoders or message passing) are frequently used
to compress structural information into a lower-dimensional latent space optimised for prediction.

% todo: handle note "word soup"...

% ---------- %

The datasets used to train such models vary in size and fidelity.

They often derive from high-throughput DFT calculations available in public repositories such as the Materials Project
or AFLOW.

% ---------- %

\subsection{What are the strengths and limitations of ML in this field?}

% ---------- %

Machine learning (ML) has emerged as a promising strategy for accelerating materials modelling, particularly in the
context of complex interfacial systems.

Its strengths lie chiefly in efficiency, scalability, and adaptability.

ML models, once trained, can rapidly approximate formation energies, surface reactivities, and stability metrics for
thousands of configurations, many of which would be prohibitively expensive to assess using first-principles methods
alone.

This computational efficiency enables exploration of broader configurational and chemical spaces, including systems with
disorder, stacking faults, or metastable structures that pose challenges for traditional simulations.

% ---------- %

Another key advantage is the capacity of ML to capture subtle, non-linear correlations in high-dimensional datasets.

When used in conjunction with tools such as ARTEMIS or RAFFLE, ML can assist not only in predicting energetics, but also
in prioritising candidate geometries during structure generation.

% ---------- %

However, ML approaches also exhibit critical limitations.

A foremost concern is their dependence on the training dataset: models typically interpolate well within the domain of
known examples but generalise poorly to unseen chemistries or extreme geometries.

This limits their reliability when extrapolating to novel material classes or interface types, unless training sets are
curated with care to ensure coverage of relevant chemical and structural motifs.

% ---------- %

Another limitation is interpretability.

While classical models (e.g. those based on density functional theory) yield physically grounded outputs, such as charge
densities or band structures, many ML models operate as statistical black boxes.

Though efforts are underway to integrate physically meaningful descriptors (e.g. n-body distribution functions, local
work function, or Bader charges), ensuring transparency and robustness remains a work in progress.

% ---------- %

Finally, ML alone cannot fully replace physics-based methods.

For systems exhibiting emergent interfacial phases, or electronic reconstructions, high-fidelity methods like DFT or
many-body perturbation theory remain necessary for validation and refinement.

As such, a hybrid strategy is often advocated; using ML to pre-screen candidate structures and flag configurations for
further high-accuracy assessment.

% ---------- %

ML constitutes a valuable addition to the materials modelling toolkit, particularly for screening
and guiding interface prediction.

Yet its performance remains contingent on data quality and physical context, and it is best deployed in tandem with more
rigorous electronic structure methods.

% ---------- %

\subsection{Machine-Learned Interatomic Potentials (MLPs)}

% ---------- %

\subsubsection{MACE}

% ---------- %

The MACE (Many-body Attention and Correlation Equivariant) model represents a recent advancement in equivariant message
passing neural networks (MPNNs), tailored for high-accuracy and scalable force field construction.

Unlike earlier MPNNs which rely primarily on two-body invariant messages, MACE employs many-body interactions and
equivariant tensor contractions to capture angular and chemical complexity with improved expressivity and data
efficiency.

% ---------- %

One notable strength of MACE is its capacity to reach state-of-the-art accuracy on diverse benchmarks, including rMD17,
3BPA, and AcAc, while requiring only two message passing iterations, due to its use of high-body-order messages.

For instance, in extrapolation tasks involving out-of-distribution molecular geometries or high-temperature
trajectories, MACE (particularly with $L=2$) has demonstrated superior performance and robustness relative to competing
models such as NequIP and BOTNet.

% ---------- %

\subsubsection{ChgNet}

% ---------- %

ChgNet is a neural network potential that incorporates both atomic positions and charge distributions into its
prediction framework, thereby explicitly learning charge transfer phenomena.

This is particularly relevant for interfaces where electrostatics, dipole formation, and local charge accumulation
significantly affect electronic properties such as band alignment.


% ---------- %

\subsubsection{Use in Interface Modelling}

% ---------- %

In the context of this project, MLPs like MACE and ChgNet offer distinct but synergistic advantages.

MACE, with its speed, high accuracy, and symmetry-respecting architecture, is well-suited for structure relaxation and
force prediction across multi-material junctions.

ChgNet, on the other hand, may provide more faithful predictions of electronic interface effects arising from local
charge redistribution.

% ---------- %

Therefore, hybrid workflows that use MLPs for broad structural exploration, followed by targeted DFT recalculations for
ambiguous configurations, are increasingly being adopted.

% ---------- %

The eventual goal is to integrate MLPs into an active learning pipeline alongside tools like RAFFLE or ARTEMIS,
facilitating the generation, screening, and optimisation of candidate interface structures at scale.

% todo: address note "should be in intro"

% ---------- %

% ---- where did this come from? ---- %
% todo find if i want to keep this!
Physics-based methods, most notably Density Functional Theory (DFT), provide high accuracy for local bonding
environments and band structure prediction.

However, the computational scaling of standard Kohn-Sham DFT---typically $\mathcal{O}(N^3)$ with respect to system
size---renders large or disordered interface simulations intractable without approximation.

Even when feasible, DFT struggles to capture extended features such as long-range strain relaxation, misfit
dislocations, and spatially diffuse defect states that dominate realistic interface behaviour.

% ---------- %

Machine learning potentials (MLPs), such as those based on graph neural networks or equivariant message passing
architectures, have emerged as promising tools to extend predictive power to larger systems at near-DFT accuracy.

Nonetheless, they too face structural limitations.

MLPs require extensive, representative training datasets, typically derived from DFT calculations, and their performance
deteriorates when applied to configurations involving unseen chemistries, surface reconstructions, or high strain
regimes.

Furthermore, many ML models lack rigorous uncertainty quantification, making it difficult to detect failure modes in
out-of-distribution predictions.

% ---------- %

Both DFT and ML approaches also fall short in capturing emergent interfacial phenomena such as charge transfer,
interface-induced dipoles, or the formation of new interfacial phases.

These effects are often strongly dependent on the local atomic registry and chemical heterogeneity, and cannot be
inferred from the properties of the isolated constituents alone.

Notably, simple models like Anderson's rule have been shown to systematically fail in 2D and 2D|3D heterostructures,
necessitating corrections such as $\Delta E_\Gamma$ and $\Delta E_\mathrm{IF}$ that explicitly depend on the atomic-scale interface
structure and electrostatics.

% ---------- %

While ML and physics-based models have complementary strengths, neither is sufficient in isolation for the reliable
prediction of interfacial behaviour.

Their limitations motivate the development of hybrid workflows, such as embedding physical constraints into ML
architectures, or coupling MLP-based screening with selective DFT refinement.

These approaches aim to achieve a balance between scalability, accuracy, and generalisability---a balance that remains at
the forefront of interface modelling challenges .

% ---^ where did this come from? ^--- %